\uuid{wGof}
\exo7id{4479}
\titre{exo7 4479}
\auteur{quercia}
\organisation{exo7}
\datecreate{2010-03-14}
\isIndication{false}
\isCorrection{false}
\chapitre{Série numérique}
\sousChapitre{Autre}
\cpge{mpsi-mp2i.raisonnement-vocabulaire-ensembliste.c, pcsi.raisonnement-vocabulaire-ensembliste.d, ptsi.raisonnement-vocabulaire-ensembliste.d, mpsi-mp2i.nombres-reels-suites.c, pcsi.nombres-reels-suites.b, ptsi.nombres-reels-suites.b, mpsi-mp2i.procedes-sommatoires-discrets.a, pcsi.series-numeriques.a, ptsi.series-numeriques.a, mpsi-mp2i.procedes-sommatoires-discrets.b, pcsi.series-numeriques.b, ptsi.series-numeriques.b, bcpst1.suites-reelles, mp-mpi.series-numeriques-vectorielles.a, pt.series-numeriques.b, bcpst2.series-reelles}
\module{Analyse}
\niveau{L1}
\difficulte{}

\contenu{
\texte{
Soit $\cal S$ l'ensemble des suites croissantes d'entiers $(q_i)$ telles que $q_0\ge 2$.
}
\begin{enumerate}
    \item \question{Si $s = (q_i) \in \cal S$, montrer que la série
$\sum_{k=0}^\infty \frac1{q_0\dots q_k}$ converge. On note $\Phi(s)$ sa somme.}
    \item \question{Montrer que l'application $\Phi : {\cal S} \to {]0,1]}$ est bijective.}
    \item \question{Soit $s = (q_i) \in \cal S$. Montrer que $\Phi(s) \in\Q$ si et seulement
    si $s$ est stationnaire.}
\end{enumerate}
}
